\section{Chapter Summary}
\label{sec:chap4_conclusion}

This chapter described the four infrastructure changes that address the highest-priority challenges identified in Chapter~3.

Section~\ref{sec:infra_improvements} replaced the BLaIR encoder with BGE-M3, introduced the NLLB-200 multilingual translation pipeline, parallelised text and CLIP encoding via an \texttt{anyio} task group, and replaced the GRU-SeqDQN recommendation module with a dual-pipeline architecture.
These changes reduce warm-cache end-to-end search latency from 1{,}102 ms to 59 ms and enable Vietnamese-language queries that previously achieved Precision@10 $\leq 0.05$ under BLaIR.

Section~\ref{sec:ingestion_improvements} replaced the batch-only Cleora rebuild with a three-stage incremental pipeline: a delta export script that queries MongoDB for new click and cart events since the last run, a Cleora augmentation step that updates the co-purchase graph, and a hot-reload step that atomically swaps the in-memory index.
The pipeline is triggered automatically at \texttt{CLEORA\_REFIT\_THRESHOLD} = 50 interactions and is fault-isolated at each stage, retaining the previous embeddings if any stage fails.

Section~\ref{sec:viz_improvements} replaced synchronous rendering with skeleton loading states, exposed the dual-pipeline recommendation architecture as two explicit sub-tabs (\textit{People Also Buy} and \textit{You Might Like}), and added an inline training visualisation showing the DIF-SASRec loss sparkline with pretrain-to-online phase annotation.
The revised interface also supports multimodal search through a combined text-and-image input area, and provides a dark mode toggle and a live recently-viewed strip.

Section~\ref{sec:storage_improvements} introduced a dual FAISS index strategy — an HNSW index for production ANN search and a flat index for exact reconstruction — loaded via memory mapping to eliminate the 45-second full-load startup penalty.
The HNSW index reduces P50 search latency from 612.19 ms to 119.56 ms (a 5.12$\times$ speedup), at a recall of 0.521 that is partially compensated by the BM25 re-ranking stage.
The Cleora behavioural embeddings are maintained as a separate hot-swappable archive covering 375{,}280 catalogue items.

Section~\ref{sec:agent_pool} described the \texttt{AgentPool} architecture that underpins concurrent personalisation: a fixed pool of eight independent DIF-SASRec instances backed by an \texttt{asyncio.Queue}, each maintaining its own weight and optimiser state, with per-user checkpoints serialised to disk between requests.
The pool provides natural backpressure under concurrency without dropped requests or shared-model race conditions, at a VRAM cost of approximately 1.18~GB for the full pool of eight agents.

Chapter~5 evaluates these changes quantitatively, reporting recommendation accuracy across 100{,}000 test users, multilingual retrieval precision across five query-type groups, and end-to-end latency across four optimisation stages.
